\section{Evaluation}
\label{sec:eval}

\subsection{Setup}
\paragraph{Subjects.}
We use three complementary sets of subjects, which Table~\ref{tab:subjects} summarizes. The first is the OWASP Benchmark for Python~\cite{owaspbench} at a pinned revision, whose 168 path-traversal test cases, 65 vulnerable and 103 safe, we split by a seeded procedure grouped by \emph{flow template}, which is the normalized code between the source and the value reaching the sink. The development set holds 34 cases in 7 families and the held-out set 134 cases in 61 families. The splitter computes families automatically from every file, the one case we read while planning went into development, and no held-out case was inspected or run through an arm before the freeze.

The second set consists of diagnostic pairs, 12 for development and 24 held out, which we wrote as small Flask handlers modelled on Python, Flask and OWASP guidance. They are author-constructed stress tests frozen before evaluation, and they cannot stand in for an independent benchmark. Family labels are disjoint between the splits, although some mechanisms such as \texttt{basename} checks and decode-after-check recur across them. The two members of a pair differ in one security-relevant way, for example a check on the value that is used against a check on some other value. A harness runs each handler under 195 combinations of traversal payloads, benign inputs, control values and input groups with the filesystem stubbed out. Every vulnerable program has a witness input whose accessed path leaves the root, and every safe program has completed runs without a witness together with a one-line static justification, although the absence of a witness among finitely many inputs is evidence that falls short of proof. The harness rejected four intended labels, which we fixed before any arm was run on held-out data, redesigning two held-out programs in the process.

The third set is a real-world case study drawn from PyVul~\cite{pyvul}. We screened its CWE-22 commits in a fixed order on feasibility alone, requiring an accessible repository, a non-test Python change and Flask request usage, and recorded each screening decision. Screening yielded six candidate pairs of vulnerable and fixed targets in four projects, which we afterwards adjudicated for applicability and mapped warning by warning to the documented defect.

\begin{table}[t]
\caption{Evaluation subjects.}
\label{tab:subjects}
\small\begin{tabular}{lrrrr}
\toprule
Corpus & Units & Vuln. & Safe & Families \\
\midrule
OWASP dev & 34 & 7 & 27 & 7 \\
OWASP test & 134 & 58 & 76 & 61 \\
Diagnostic dev & 24 & 12 & 12 & 12 \\
Diagnostic held-out & 48 & 24 & 24 & 12 \\
\bottomrule
\end{tabular}

\end{table}

\paragraph{Protocol.}
Before any arm was run on test data, we froze and tagged in version control the Semgrep rule and its hash, the split, the diagnostic manifest, the p2 prompts, the analyzer, the gate and the scorer. The scoring unit is the test case or program, which counts as reported if any of its alerts is retained. We report precision, recall, F1, TP retention $|R\cap G|/|A\cap G|$ and FP removal $|(A\setminus R)\setminus G|/|A\setminus G|$, where $A$ is the set of baseline-reported units, $R$ the units an arm retains and $G$ the vulnerable units. Uncertainty is a 95\% percentile bootstrap over families with 2{,}000 resamples. The models are \texttt{qwen3-coder-next} as the primary model and \texttt{qwen3.5:9b} with thinking turned off, both served by SoCLaaS at temperature 0 with JSON output.

\subsection{RQ1: The removal/retention trade-off}
Table~\ref{tab:main} gives the main results, where TP retention is the share of baseline true positives an arm keeps and FP removal the share of baseline false positives it drops. On OWASP the baseline reports 68 of the 134 held-out cases, of which 35 are vulnerable and 33 safe, and its 23 misses all come from two engine limitations of Semgrep CE. Fifteen involve taint passing through string slices and eight involve a value stored in a dict by subscript and read back by key, and since every arm only filters baseline warnings, no refinement can recover them.

D removes 12 of the 33 false positives without losing a vulnerability, and in all 12 the attacker value reaches the sink only through a branch or conditional expression whose condition folds to a constant. The 21 false positives it keeps are fully analysed but tainted. Eight are protected only by a block-list such as \texttt{'../' in bar}, and although their official label is safe we mark them \emph{ambiguous}, because a bare \texttt{..} still reaches the parent directory. Five rely on a prefix check that compares an absolute path with a relative string and so always rejects, and four rely on \texttt{ConfigParser} key separation, in which the tainted value is stored under one key and a different key is read, which our analysis does not model. The last four take the attacker value from the first path segment, which the benchmark's routing always fixes to the constant \texttt{benchmark}, so they are safe for a reason that no request-data threat model captures. All scores use the official labels, but leaving out the eight ambiguous cases raises the FP removal of D and E from 0.36 to 0.48 and lowers that of the small model's L from 0.70 to 0.64, since seven of its removals are ambiguous cases. Part of the apparent advantage of that arm therefore rests on disputed labels.

With both models E retains the same OWASP units as D, and we found no OWASP case in which structured evidence let the model go beyond the recognizers. L with the primary model removes 11 false positives but wrongly dismisses one real vulnerability.

The held-out diagnostic pairs are built to contain non-standard guards, and on them the differences between arms grow larger. D removes 13 of 19 false positives with full retention, while L removes 15 at the cost of 7 of 23 real vulnerabilities. E and EC remove the same 15 as L and lose only 2.

\paragraph{Uncertainty.}
The intervals are wide, and D's FP removal on OWASP, for example, spans [0.00, 0.67], so a difference of one or two cases between arms says little about behaviour in general. A TP-retention interval of [1.00, 1.00] records only that no loss occurred in the sample and places no bound on future losses, and the diagnostic intervals describe variation within our 12 constructed families alone. The artifact tables give the family-bootstrap interval of every cell in Table~\ref{tab:main}. We also bootstrap paired differences between arms on the same resamples. On the held-out pairs, structured evidence reduces wrongful suppressions relative to ordinary prompting, with E$-$L at $-5$ [$-10$, $-1$] for \texttt{qwen3-coder-next} and $-8$ [$-14$, $-3$] for \texttt{qwen3.5:9b}, whereas on OWASP every paired difference between arms has an interval that includes zero.

\begin{table*}[t]
\caption{Held-out results at unit level.}
\label{tab:main}
\small\begin{tabular}{llrrrrrrrr}
\toprule
Subjects / model & Arm & TP & FP & FN & Prec. & Rec. & F1 & TP ret. & FP rem. \\
\midrule
\multirow{5}{*}{\shortstack[l]{OWASP test\\\texttt{qwen3-coder-next}}} & S & 35 & 33 & 23 & 0.515 & 0.603 & 0.556 & 1.000 & 0.000 \\
 & D & 35 & 21 & 23 & 0.625 & 0.603 & 0.614 & 1.000 & 0.364 \\
 & L & 34 & 22 & 24 & 0.607 & 0.586 & 0.596 & 0.971 & 0.333 \\
 & E & 35 & 21 & 23 & 0.625 & 0.603 & 0.614 & 1.000 & 0.364 \\
 & EC & 35 & 22 & 23 & 0.614 & 0.603 & 0.609 & 1.000 & 0.333 \\
\midrule
\multirow{5}{*}{\shortstack[l]{OWASP test\\\texttt{qwen3.5:9b}}} & S & 35 & 33 & 23 & 0.515 & 0.603 & 0.556 & 1.000 & 0.000 \\
 & D & 35 & 21 & 23 & 0.625 & 0.603 & 0.614 & 1.000 & 0.364 \\
 & L & 29 & 10 & 29 & 0.744 & 0.500 & 0.598 & 0.829 & 0.697 \\
 & E & 35 & 21 & 23 & 0.625 & 0.603 & 0.614 & 1.000 & 0.364 \\
 & EC & 35 & 26 & 23 & 0.574 & 0.603 & 0.588 & 1.000 & 0.212 \\
\midrule
\multirow{5}{*}{\shortstack[l]{Diagnostic held-out\\\texttt{qwen3-coder-next}}} & S & 23 & 19 & 1 & 0.548 & 0.958 & 0.697 & 1.000 & 0.000 \\
 & D & 23 & 6 & 1 & 0.793 & 0.958 & 0.868 & 1.000 & 0.684 \\
 & L & 16 & 4 & 8 & 0.800 & 0.667 & 0.727 & 0.696 & 0.789 \\
 & E & 21 & 4 & 3 & 0.840 & 0.875 & 0.857 & 0.913 & 0.789 \\
 & EC & 21 & 4 & 3 & 0.840 & 0.875 & 0.857 & 0.913 & 0.789 \\
\midrule
\multirow{5}{*}{\shortstack[l]{Diagnostic held-out\\\texttt{qwen3.5:9b}}} & S & 23 & 19 & 1 & 0.548 & 0.958 & 0.697 & 1.000 & 0.000 \\
 & D & 23 & 6 & 1 & 0.793 & 0.958 & 0.868 & 1.000 & 0.684 \\
 & L & 15 & 4 & 9 & 0.789 & 0.625 & 0.698 & 0.652 & 0.789 \\
 & E & 23 & 7 & 1 & 0.767 & 0.958 & 0.852 & 1.000 & 0.632 \\
 & EC & 23 & 7 & 1 & 0.767 & 0.958 & 0.852 & 1.000 & 0.632 \\
\bottomrule
\end{tabular}

\end{table*}

\subsection{RQ2: Do evidence and the gate help?}

\paragraph{Evidence versus ordinary prompting.}
With the same source and budget, the evidence-guided configuration reduced wrongful suppressions on the held-out pairs from 7 to 2 at equal FP removal, and on OWASP it turned L's 11 removals and one loss into 12 removals without a loss. L's seven mistakes on the pairs include both decode-after-check programs, both prefix-confusion programs, a guard that only logs and a reachable \texttt{match} default branch. In each of them the model judged a nearby check or a reassuring name when the value that mattered was the one consumed at the sink, and the evidence record states that fact explicitly. E differs from L in its representation, its obligations and the analysis conclusions it is given, and the comparison cannot separate the effects of these three changes. One of E's retentions, on a program without an evidence record, and one of L's, on an invalid reply, are conservative fallbacks that involve no model judgment.

\paragraph{The gate.}
Table~\ref{tab:gate} sorts E's suppressions by gate category, and a rejection counts as a benefit when the alert is vulnerable and as a cost when it is safe. The gate never rejected a wrongful E suppression on held-out data, while on OWASP it rejected one \emph{correct} suppression with the primary model and five with the smaller one. With the primary model the rejected case was safe because its value is constant, although E justified the suppression with a guard that does not dominate the sink, so EC removed one false positive fewer than E. Every EC suppression on OWASP falls in the \texttt{structural\_check\_passed} category, so on that corpus EC can at best equal D.

On the held-out pairs EC goes beyond D in four programs, all labelled \texttt{semantic\_judgment\_remaining}. Two of these suppressions are correct, covering canonicalization followed by a check and \texttt{basename} combined with rejection of \texttt{..}, and two are wrong, covering \texttt{startswith(ROOT)} without a separator and a containment check that the attacker can switch off with a second parameter. The gate verifies that a guard dominates the sink and fits the dependency structure, but it cannot verify that the predicate confines the path, and both wrong suppressions fail on that point. In these runs the gate improved no measured retention and cost some precision.

\paragraph{Versus deterministic refinement.}
Against deterministic refinement the LLM arms gained little. D kept every benchmark-labelled true positive on both held-out corpora, and on OWASP it removed 12 of 33 false positives against 11 for EC with the primary model and 7 with the small one. On the pairs it removed 13 of 19 against 15 for E and EC with the primary model, and those two extra removals cost two real vulnerabilities.

\begin{table}[t]
\caption{E's suppressions under the EC gate.}
\label{tab:gate}
\resizebox{\columnwidth}{!}{\begin{tabular}{lrrrrrr}
\toprule
Corpus, model & E supp. & Struct. & Contract & Semantic & \multicolumn{2}{c}{Rejected: wrong / correct} \\
\midrule
OWASP, coder-next & 12 & 11 & 0 & 0 & 0 & 1 \\
OWASP, 9b & 12 & 7 & 0 & 0 & 0 & 5 \\
Diagnostic, coder-next & 17 & 3 & 10 & 4 & 0 & 0 \\
Diagnostic, 9b & 12 & 2 & 10 & 0 & 0 & 0 \\
\bottomrule
\end{tabular}
}
\end{table}

\subsection{RQ3: Patterns}
Table~\ref{tab:pairs} reports per family how many held-out pairs each arm solves, counting a pair as solved when both members are classified correctly, and its last column counts the wrongful suppressions of L, E and EC, while Table~\ref{tab:cases} walks through representative cases. D, E and EC each solve 17 of the 24 pairs and L solves 13, although D and E succeed on \emph{different} families.

D solves every pair whose safety it can prove, which covers overwrites, unreachable flows, \texttt{realpath} and \texttt{abspath} containment, and checks on the wrong value or exits that only log. It fails where safety rests on a block-list or on \texttt{basename}, where containment follows \texttt{normpath}, which its recognizer does not accept, or where two definitions reach the sink. E and EC recover the pairs built on \texttt{normpath} canonicalization and on \texttt{basename} with rejection, but they lose the prefix-confusion and attacker-switchable guard pairs that D gets right.

Two pairs defeat every arm for reasons outside triage. In the first named-sanitizer pair the helper body lives in another function, so no arm sees it, and in the first aliasing pair Semgrep misses the vulnerable member because its \texttt{append} propagator taints the alias and leaves the list itself untainted, which is the baseline's only false negative on the pairs.

\begin{table}[t]
\caption{Held-out diagnostic pair accuracy and wrongful suppressions.}
\label{tab:pairs}
\resizebox{\columnwidth}{!}{\begin{tabular}{lrrrrrr}
\toprule
Family & S & D & L & E & EC & Wrongful L/E/EC \\
\midrule
absolute\_join & 0/2 & 0/2 & 1/2 & 0/2 & 0/2 & 0/0/0 \\
alternate\_branch & 0/2 & 1/2 & 1/2 & 0/2 & 0/2 & 1/1/1 \\
canonicalize\_only & 0/2 & 1/2 & 2/2 & 2/2 & 2/2 & 0/0/0 \\
checked\_vs\_original & 0/2 & 2/2 & 0/2 & 2/2 & 2/2 & 2/0/0 \\
container\_alias & 1/2 & 1/2 & 1/2 & 1/2 & 1/2 & 0/0/0 \\
guard\_other\_value & 0/2 & 2/2 & 2/2 & 2/2 & 2/2 & 0/0/0 \\
named\_sanitizer & 0/2 & 0/2 & 1/2 & 1/2 & 1/2 & 0/0/0 \\
overwrite & 1/2 & 2/2 & 2/2 & 2/2 & 2/2 & 0/0/0 \\
prefix\_confusion & 0/2 & 2/2 & 0/2 & 1/2 & 1/2 & 2/1/1 \\
reject\_exits & 0/2 & 2/2 & 1/2 & 2/2 & 2/2 & 1/0/0 \\
send\_from\_directory\_base & 2/2 & 2/2 & 2/2 & 2/2 & 2/2 & 0/0/0 \\
unreachable\_flow & 0/2 & 2/2 & 0/2 & 2/2 & 2/2 & 1/0/0 \\
\midrule
Total & 4/24 & 17/24 & 13/24 & 17/24 & 17/24 & \\
\bottomrule
\end{tabular}
}
\end{table}

\begin{table*}[t]
\caption{Representative cases with the primary model, where k means kept and s suppressed.}
\label{tab:cases}
\small
\begin{tabular}{lcccccl}
\toprule
Case (truth) & D & L & E & EC & Gate & What happened \\
\midrule
checked\_vs\_original1 (vuln) & k & s & k & k & -- & L trusts the check on \texttt{candidate}, E sees \texttt{open} use the decoded \texttt{name} \\
unreachable\_flow2 (vuln) & k & s & k & k & -- & L calls the reachable \texttt{match} default rare \\
canonicalize\_only2 (safe) & k & s & s & s & semantic & \texttt{normpath} with \texttt{startswith(ROOT+sep)}, correct beyond D \\
prefix\_confusion1 (vuln) & k & s & s & s & semantic & \texttt{startswith(ROOT)} wrongly accepted as containment \\
OWASP 00091 (safe) & s & s & s & k & rejected & E labels a constant value as guarded \\
\bottomrule
\end{tabular}
\end{table*}

\paragraph{Real-world case study.}
For the real-world targets we read every fix diff after the runs and mapped each warning to the documented defect without looking at arm decisions, storing the mapping in \texttt{data/realworld/*\_adjudication.csv}. Two of the six candidate targets fall outside the task, one being a CORS path-matching flaw with no filesystem operation and the other a Windows drive-letter check. Among the four applicable targets, whoogle-search is the clearest, since S reports the documented flow, in which a request name is joined to the configuration directory and opened, through three warnings that every arm keeps. A fourth warning in its changed functions is \texttt{send\_file(io.BytesIO())}, which takes no path and belongs to an SSRF fix in the same commit, so we count it as unrelated. In mlflow \texttt{5044878da0}, S reports two warnings on the raw request path that the fix replaces with its validated form. Whether these local sinks were exploitable before the fix remains open, so we count them only as fix-related, and D, E and EC keep both while L drops both. For opendiamond and mlflow \texttt{ac4b697bb0} S reports nothing, because the attacker value is a route variable, which is not a declared source, or the filesystem operation lies in another function.

At the fixed revisions the defect sites are still reported, since the fixes rely on an allow-list regex and a validator that Semgrep CE does not model, and D, E and EC keep all five such warnings whereas L keeps one. With so few changed functions, the case study illustrates behaviour on real code and gives no estimate of real-world detection or retention rates.

\subsection{RQ4: Stability, model size and cost}
\paragraph{Fresh generations.}
We drew 30 held-out alerts by a seeded rule, stratified by E's frozen decision (19 OWASP alerts, 11 diagnostic). For each, L and E were re-run three times with new replicate IDs, giving 180 genuinely fresh calls at temperature 0.
\begin{itemize}
\item \textbf{L} repeated its frozen decision in all 90 calls.
\item \textbf{E} changed on one alert, \texttt{prefix\_confusion2} (vulnerable). The frozen response kept it, but 2 of 3 fresh generations suppressed it, and the gate accepted them as \texttt{semantic\_judgment\_remaining}.
\end{itemize}
Decoding at temperature 0 is therefore not deterministic. E's held-out count of two wrongful suppressions could as well have been three, and that variation falls exactly in the category the gate cannot check.

\paragraph{Smaller model.}
With \texttt{qwen3.5:9b} (thinking disabled), ordinary triage became more aggressive and more harmful. On OWASP, L removed 23 of 33 false positives but discarded 6 of 35 real vulnerabilities (TP retention 0.829). On the pairs, it lost 8 of 23 (0.652). E kept every real vulnerability on both corpora, two of them on the pairs only by fallback (no evidence record, and one invalid reply). It added nothing over D: identical on OWASP, and one removal fewer on the pairs (0.632 versus 0.684).

The gate behaved differently with this model. It rejected five correct E suppressions on OWASP, so EC's FP removal fell from 0.364 to 0.212. In every case the model cited line numbers instead of the evidence IDs it was asked for. The gate's value thus depends on the model following the citation protocol. With a weaker model, its rejections here cost precision and recovered no vulnerability.

\paragraph{Cost.}
Table~\ref{tab:resources} lists the LLM cost of each held-out run.
\begin{itemize}
\item \textbf{Speed.} The primary model answered in a median of 1.9\,s per cold call (90th percentile 3.2\,s), the small model in 13.0\,s (32.8\,s). Each LLM arm over the 68 OWASP alerts took 2.5 to 13 minutes.
\item \textbf{Rate limiting.} SoCLaaS limited bursts: 171 transport retries were logged across the four held-out runs. Transport retries did not depend on answer content, and the only content-dependent retry was the declared schema repair (6 requests in total).
\item \textbf{Total use.} The whole project, development and the post-review correction runs included, used 874 uncached requests and 2.0M tokens, within the planned caps of 1{,}500 and 12M. Table~\ref{tab:resources} counts completion requests, which include schema repairs: 68 OWASP alerts give 136 triage jobs but 137 requests.
\item \textbf{Deterministic parts.} Semgrep and evidence extraction took under 2\,s per corpus.
\item \textbf{Price.} SoCLaaS has no charge or quota for SoC users. The figures matter as shared compute, and because they bound what a larger study would need.
\end{itemize}


\begin{table}[t]
\caption{LLM resources per held-out run.}
\label{tab:resources}
\resizebox{\columnwidth}{!}{\begin{tabular}{llrrrrr}
\toprule
Run & Model & Requests & Uncached & In tok. & Out tok. & Time (s) \\
\midrule
owasp-qcn & \texttt{qwen3-coder-next} & 137 & 137 & 322,782 & 27,529 & 347 \\
owasp-q9b & \texttt{qwen3.5:9b} & 136 & 136 & 351,325 & 21,446 & 1529 \\
diag-qcn & \texttt{qwen3-coder-next} & 87 & 87 & 104,463 & 27,143 & 235 \\
diag-q9b & \texttt{qwen3.5:9b} & 84 & 84 & 99,690 & 13,835 & 500 \\
\bottomrule
\end{tabular}
}
\end{table}

\subsection{Post-hoc extension}
After seeing the two wrongful EC suppressions we tried a variant, \textbf{EC+}, which adds two checks before deferring to the model on a guard. It rejects containment-shaped guards that fail the sound recognizer as well as guards whose truth depends on other attacker input. Replaying the frozen E responses without new model calls, EC+ reaches TP retention 1.000 and FP removal 0.737 on the held-out pairs, against 0.913 and 0.789 for EC, and it leaves the OWASP test results unchanged. On OWASP development data it \emph{loses} three correct suppressions in which E cited an unsound prefix guard for values that were constant. Since EC+ was designed after inspecting these failures, its numbers are exploratory and do not count as held-out results.

\subsection{Threats to validity}
The OWASP Python benchmark is a preliminary release whose cases are generated from templates, and we score against its labels as given, even though five of its ``safe'' cases are safe only because their check always fails at runtime and eight block-list cases are ambiguous, as the sensitivity analysis in RQ1 shows. We wrote the diagnostic programs ourselves, so their families reflect our view of what is hard, and their labels rest on finite execution witnesses plus short static arguments. The review in Section~\ref{sec:correction} found real soundness bugs in the frozen extractor, and although the corrected version changes no held-out result, bugs of other kinds may remain.

Samples are small and some error types have few clusters. The two models also differ in family, size, serving stack and reasoning mode, so the model comparison cannot isolate any one of these factors. SoCLaaS may change the model behind a name, which we mitigate by recording the served model and backend fingerprint of every call, and all calls to the primary model were served under one fingerprint. Finally, public benchmark code may appear in the training data of the models.
